\section{The retained-label actual coefficient interface}
\label{sec:tpc52-interface}

We reconstruct the operator before proving an estimate.  The order of
labels and norms is part of the theorem: a diagonal Mellin identity is
available only before coherent prime summation and residual grouping.
The interface is inherited from TPC--25, TPC--36, and TPC--47--51
\citep{WangTPC25,WangTPC36,WangTPC47,WangTPC48,WangTPC50,WangTPC51}.

\subsection{Rows, the literal mask, and the actual profile}

Fix the physical shift \(h_0\ne0\) and one source-compatible dyadic
packet.  A row is
\begin{equation}
 \rho=(\ell_\rho,d_\rho),
 \qquad m_\rho=\ell_\rho d_\rho\asymp Q,
 \label{eq:tpc52-row}
\end{equation}
where \(\ell_\rho\) is prime, \(d_\rho\) is squarefree, and the inherited
size separation makes \(m_\rho\) determine \(\rho\).  Let
\(\cA_1,\cA_2\) be the source and opposite-row systems.  The literal
static mask is
\begin{equation}
 K_{\alpha\gamma}
 :=\one_{\ell_\alpha\ne\ell_\gamma}
   \one_{|m_\alpha-m_\gamma|>\Delta}
   \one_{(d_\alpha,d_\gamma)\le G},
 \qquad \Delta=QX^{-\kappa_0},\quad G=X^{\kappa_0}.
 \label{eq:tpc52-literal-mask}
\end{equation}
It is not replaced by a signed rank-one representation in this section.

Before the mask is decomposed, the joint multiplier is
\begin{equation}
 \mathfrak A_{m,\gamma}^{\rm act}(j)
 =K_{m\gamma}\Xi_{m,\gamma}(j)
  \mathcal W_{m,\gamma}(j/J),
 \label{eq:tpc52-joint-multiplier}
\end{equation}
where \(\Xi\) has fixed period \(H_0=O_{h_0}(1)\), and
\(\mathcal W\) is the complete unseparated product of the declared
smooth row--orbit factors.  On the left, the literal source coefficient
has the form
\begin{equation}
 \gamma_m^{(1)}
 =\mu(d(m))(\log\ell(m))
  \omega_D(d(m))\psi_L(\ell(m)/L)\zeta_m^{(1)}.
 \label{eq:tpc52-source-coefficient}
\end{equation}
The factor \(\zeta_m^{(1)}\) is kept exactly; no lower bound or source
equidistribution is assumed below.

Fix one complete source/opposite-row residue group of positive density and
suppress its outer label.  A fixed finite collection of such row groups is
handled later by orthogonal summation.  Let \(\cA_0\) be the direct orbit
residue cells inside the chosen row group, let \(\cO_L\) be its common
opposite-row set, and let \(\cI_X\) contain the supported pairs \((m,r)\)
before prime summation.  We identify \(a\in\cA_0\) with its orbit
residue modulo \(H_0\) and put
\begin{equation}
 \begin{aligned}
 \cV_L&:=\ell^2(\cO_L\times\cS_X),\\
 \cJ_{X,a}&:=\{j\in\Z:j/J\in I_0,
                    \ j\equiv a\!\!\pmod {H_0}\},\\
 \cH_X^{\rm coeff}
 &:=\bigoplus_{a\in\cA_0}\ \bigoplus_{(m,r)\in\cI_X}
      \ell^2(\cJ_{X,a};\cV_L).
 \end{aligned}
 \label{eq:tpc52-coefficient-space}
\end{equation}
where \(I_0\Subset(0,\infty)\) is a fixed common normalized orbit
interval.  All index sets are finite for each \(X\).

For a direct residue cell \(a\in\cA_0\), the actual profile is
\begin{equation}
 G_{m,r,a}^{L}(j)
 :=\sum_{\gamma\in\cO_L}
   \gamma_m^{(1)}K_{m\gamma}
   \Xi_{m,\gamma}(j)\mathcal W_{m,\gamma}(j/J)
   e_{\gamma,d(m)r}\in\cV_L,
 \qquad j\in\cJ_{X,a}.
 \label{eq:tpc52-actual-vector-profile}
\end{equation}
Thus the coefficientwise physical vector is
\begin{equation}
 G_X^L:=\bigoplus_{a\in\cA_0}\ \bigoplus_{(m,r)\in\cI_X}
             (G_{m,r,a}^{L}(j))_{j\in\cJ_{X,a}}
 \in\cH_X^{\rm coeff}.
 \label{eq:tpc52-actual-coefficient-vector}
\end{equation}
Because all displayed labels are orthogonal,
\begin{align}
 \|G_X^L\|_{\cH_X^{\rm coeff}}^2
 &={}
 \sum_{a\in\cA_0}\sum_{(m,r)\in\cI_X}
 \sum_{j\in\cJ_{X,a}}
 \sum_{\gamma\in\cO_L}
 |\gamma_m^{(1)}|^2K_{m\gamma}
 \notag\\[-2pt]
 &\hspace{5em}\times
 |\Xi_{m,\gamma}(j)|^2
 |\mathcal W_{m,\gamma}(j/J)|^2.
 \label{eq:tpc52-exact-physical-energy}
\end{align}
There is no cross term between different \((m,r,j,\gamma)\).  In
particular, two pairs \((m,r)\) that later produce the same residual
label have not yet been grouped.

\subsection{The canonical Mellin synthesis}

Direct residue splitting and real-axis Mellin inversion give a fixed
complex provenance measure \(\lambda_X\) on a finite disjoint union of
spaces \(F_a\times\R^{q_a}\).  Write its polar decomposition as
\begin{equation}
 d\lambda_X(\theta)=\varsigma_X(\theta)d\mu_X(\theta),
 \qquad \mu_X=|\lambda_X|,
 \qquad |\varsigma_X|=1\quad\mu_X\text{-a.e.}
 \label{eq:tpc52-polar-measure}
\end{equation}
For a direct cell, its literal-mask layer
\(k_{X,\theta}^L\in\cH_X^{\rm coeff}\) has coordinates
\begin{align}
 k_{X,\theta}^L(a;m,r,j)
 :={}&u_{a,\tau}^L(m)\phi_{a,\tau}(j/J)
 \notag\\
 &\times\sum_{\gamma\in\cO_L}
 K_{m\gamma}v_{a,\tau}(\gamma)e_{\gamma,d(m)r}.
 \label{eq:tpc52-canonical-layer}
\end{align}
No projective-mask factor \(p_\eta q_\eta\) occurs in
\eqref{eq:tpc52-canonical-layer}.  Define
\begin{equation}
 A_{X,\theta}^L:=\varsigma_X(\theta)k_{X,\theta}^L,
 \qquad
 \cS_X^Lf:=\int f(\theta)A_{X,\theta}^L\,d\mu_X(\theta).
 \label{eq:tpc52-canonical-synthesis}
\end{equation}
The transform densities are Schwartz, while the finite-dimensional layer
norm has at most polynomial growth in the Mellin variables.  Hence its norm
is integrable against the transform variation, and
\eqref{eq:tpc52-canonical-synthesis} is a genuine Bochner integral on its
natural domain.

The provenance-fixed direction is \(f=\one\).  Coordinatewise Mellin
inversion gives the exact identity
\begin{equation}
 \boxed{\cS_X^L\one=G_X^L.}
 \label{eq:tpc52-actual-direction}
\end{equation}
For a fixed polynomial profile weight \(w_B(\theta)\ge1\), the displayed
canonical envelope is
\begin{equation}
 \cE_{1,B}^L(f)
 :=\int w_B(\theta)|f(\theta)|
 \|A_{X,\theta}^L\|_{\cH_X^{\rm coeff}}\,d\mu_X(\theta),
 \label{eq:tpc52-canonical-envelope}
\end{equation}
and the actual-direction condition number is
\begin{equation}
 \kappa_{\rm can,B}^L(X)
 :=\frac{\cE_{1,B}^L(\one)^2}
         {\|G_X^L\|_{\cH_X^{\rm coeff}}^2},
 \label{eq:tpc52-canonical-condition}
\end{equation}
when \(G_X^L\ne0\).

\subsection{Normalized profile geometry}

The original provenance contains fixed smooth factors such as
\(W(m_\gamma j/X)\) and \(\Psi(d_\gamma j/K)\)
\citep{WangTPC25}.  Put
\begin{equation}
 x_\gamma=\log(\ell_\gamma/L),
 \qquad y_\gamma=\log(d_\gamma/D),
 \qquad t_j=j/J.
 \label{eq:tpc52-log-coordinates}
\end{equation}
Then
\begin{equation}
 \frac{m_\gamma j}{X}
 =\frac{LDJ}{X}\e^{x_\gamma+y_\gamma}t_j,
 \qquad
 \frac{d_\gamma j}{K}
 =\frac{DJ}{K}\e^{y_\gamma}t_j.
 \label{eq:tpc52-scaled-arguments}
\end{equation}
The packet relations \(LDJ\asymp X\) and \(DJ\asymp K\) keep the two
leading scale factors in fixed compact subsets of \((0,\infty)\).

After one of finitely many direct residue and smooth subdivisions, we
write the modulus of the unseparated profile as
\begin{equation}
 F_\beta(x,y,t)
 =R_\beta(x,y,t)
  \prod_{\nu=1}^q
  W_\nu(z_{\nu,\beta}(x,y,t)),
 \qquad \beta\in\mathfrak B,
 \label{eq:tpc52-profile-factorization}
\end{equation}
where:
\begin{enumerate}[label=\textup{(\alph*)},leftmargin=2.5em]
 \item \(\mathfrak B\) is a fixed compact set of normalized source-scale
 parameters;
 \item \(R_\beta\) is the carrier formed from the remaining fixed
 cutoffs and the direct residue amplitude;
 \item each \(W_\nu\in C_c^\infty((0,\infty))\) is a declared physical
 shape; and
 \item the maps \(z_{\nu,\beta}\) and all displayed profiles form a
 uniformly bounded \(C^\infty\) family on
 \begin{equation}
  \Omega:=[0,\log2]^2\times I_0.
  \label{eq:tpc52-profile-domain}
 \end{equation}
\end{enumerate}
Typical arguments are
\begin{equation}
 z_{1,\beta}(x,y,t)=c_\beta\e^{x+y}t,
 \qquad
 z_{2,\beta}(x,y,t)=c'_\beta\e^yt,
 \label{eq:tpc52-typical-profile-arguments}
\end{equation}
with \(c_\beta,c'_\beta\) in compact positive intervals.  Formula
\eqref{eq:tpc52-profile-factorization} is an assumption on the declared
direct packet, not an assertion that the two examples in
\eqref{eq:tpc52-typical-profile-arguments} exhaust every inherited
smooth factor.

\subsection{The downstream boundary}

Only after \eqref{eq:tpc52-actual-coefficient-vector} is formed does the
terminal map introduce primes and the determinant slice:
\begin{equation}
 \cF_{h_0}^{\rm act,L}
 =\sum_{\substack{a,p,r,m,j:\,pr-mj=h_0\\
                    (m,r)\in\cI_X,\ j\in\cJ_{X,a}}}
 \mu(r)\Lambda(pr)G_{m,r,a}^L(j)\otimes e_j.
 \label{eq:tpc52-downstream-prime-map}
\end{equation}
At fixed \((m,r)\), the \(p\)-sum is coherent.  The present paper neither
moves \eqref{eq:tpc52-canonical-envelope} through this map nor reverses
the resulting Minkowski inequality.  Likewise, the maps grouping
\(s=d(m)r\), completing frozen shifts, and identifying the native affine
Gram all occur downstream.
